% TOP_PROBLEM: {"id": 1200010, "title": "Some Questions — Question 10", "category_id": 17, "category": {"id": 17, "name": "group_theory", "display_name": "Group Theory", "description": "Problems about groups, group actions, representations, and related algebraic structures.", "slug": "group-theory", "order_index": 17, "created_at": "2026-07-31T15:26:25.671Z"}, "status": "open", "problem_number": "AMR-011-0010", "set_id": 13, "statusQualification": "Makes infinitude explicit to exclude the immediate finite-group counterexample and preserve the source expander-family setting."}
% FIRST_POSTED: 2026-10-01
% ENTRY_KIND: statement_only
% UnsolvedMath ID 1200010; source catalogue code AMR-011-0010
% !TeX program = lualatex
% Definition and sources only; this file is not an LLM solution attempt.
\documentclass[11pt,a4paper]{article}
\usepackage[margin=25mm]{geometry}
\usepackage{fontspec}
\setmainfont{lmroman10-regular.otf}[BoldFont=lmroman10-bold.otf,
 ItalicFont=lmroman10-italic.otf,BoldItalicFont=lmroman10-bolditalic.otf]
\usepackage{amsmath,amssymb,mathtools}
\usepackage{unicode-math}
\setmathfont{latinmodern-math.otf}
\AtBeginDocument{\let\setminus\smallsetminus}
\usepackage{xurl,hyperref}
\hypersetup{colorlinks=true,urlcolor=blue}
\setcounter{secnumdepth}{0}
\setlength{\parindent}{0pt}
\setlength{\parskip}{5pt}
\setlength{\emergencystretch}{3em}
\begin{document}
\section*{Some Questions — Question 10}\label{1200010}
{\small UnsolvedMath ID 1200010 · AMR-011-0010 · Group Theory · AMR Open Problem Lists}
\subsection{Definitions and mathematical statement}
Let $\Gamma$ be an infinite finitely presented group, and let
$\Gamma=\Gamma_0\geq\Gamma_1\geq\Gamma_2\geq\cdots$ be finite-index normal subgroups with $\bigcap_n\Gamma_n=\{1\}$. Choose a finite symmetric generating set $S$. On functions on $\Gamma/\Gamma_n$, let
\[
 P_n f(q)=|S|^{-1}\sum_{s\in S}f(qs).
\]
Assume property $\tau$ along the chain: there is $\varepsilon>0$ such that $\langle(I-P_n)f,f\rangle\geq\varepsilon\|f\|_2^2$ for every mean-zero function and every $n$, using uniform measure on each quotient. This specifies uniform expansion of the finite quotient Cayley graphs.

Must $\Gamma$ contain a subgroup isomorphic to the free group $F_2$? Equivalently, must there exist two elements such that every nonempty freely reduced word in them and their inverses is nonidentity?

Finite presentation, normality, trivial intersection, and the uniform spectral gap are simultaneous hypotheses. Neither merely nonamenability nor a collection of finite quotients without a separating chain is the stated premise. The free subgroup need not have finite index or be detected as a free group inside any individual finite quotient. Infinitude is made explicit: a finite group admits an eventually trivial chain and vacuous uniform finite-quotient bounds but cannot contain $F_2$. For an infinite group with trivial intersection, the quotient orders tend to infinity, which is the expander-family setting of the source discussion.
\subsection{Short English statement}
Does an infinite finitely presented group with a separating expanding chain of finite quotients contain a nonabelian free subgroup?
\subsection{Sources}
\begin{itemize}
\item[S1] ULAM.AI and the UnsolvedMath Contributors, supplied problems.json, record 1200010 (AMR-011-0010), AMR Open Problem Lists. Catalogue identity and source transcription; CC BY 4.0.
\url{https://huggingface.co/datasets/ulamai/UnsolvedMath}.
\item[S2] Abert - Some questions (pdf) (2010), Question 10. Original source indexed by AMR-011-0010.
\url{https://www.renyi.hu/~abert/questions.pdf}.
\end{itemize}
\end{document}
